\chapter{Réalisation}
\label{chap:realisation}

\begin{introchapitre}
Ce chapitre présente la mise en œuvre de SentinelOps : l'environnement de travail, les choix
technologiques, les extraits de code les plus importants, le lanceur en une commande et les
interfaces de l'application.
\end{introchapitre}

\section{Environnement de travail}

Le développement a été réalisé sur un poste \afaire{marque, processeur, mémoire} sous Windows avec Docker Desktop (WSL~2), et les tests de déploiement sur une instance cloud Linux (Ubuntu) \afaire{fournisseur et taille de l'instance}. Le tableau~\ref{tab:logiciel} résume l'environnement logiciel.


\begin{table}[htbp]
\centering
\small
\begin{tabularx}{\linewidth}{L{3.6cm}L{2.4cm}Y}
\toprule
\textbf{Outil} & \textbf{Version} & \textbf{Utilisation} \\
\midrule
Python & 3.11 & agent de supervision (asyncio, psutil, docker, PyYAML) \\
Node.js & 24 & serveur du dashboard (Express, ws, better-sqlite3) \\
React, TypeScript & 18, 5 & interface utilisateur (Vite, Recharts) \\
SQLite & 3 (mode WAL) & stockage partagé \\
Docker, Docker Compose & 27+, v2 & conteneurisation et orchestration locale \\
docker-socket-proxy & 0.5.0 & accès restreint à l'API Docker \\
pytest, node:test & --- & tests automatisés \\
GitHub Actions & --- & intégration continue \\
Trivy & --- & analyse de vulnérabilités des images \\
Visual Studio Code, Git & --- & édition du code et gestion de versions \\
\bottomrule
\end{tabularx}
\caption{Environnement logiciel}
\label{tab:logiciel}
\end{table}

\section{Choix technologiques}

\begin{table}[htbp]
\centering
\small
\begin{tabularx}{\linewidth}{L{3cm}L{3.4cm}Y}
\toprule
\textbf{Besoin} & \textbf{Choix} & \textbf{Justification} \\
\midrule
Agent de collecte & Python + asyncio & bibliothèques mûres pour le système (\code{psutil}) et Docker ; une boucle asynchrone suffit pour une cadence de 5 s \\
Serveur web & Node.js + Express & un seul langage avec l'interface ; WebSocket natif via \code{ws} \\
Interface & React + TypeScript & composants réutilisables, typage des données de l'API \\
Stockage & SQLite (WAL) & aucun serveur à installer, un fichier sur un volume, lectures concurrentes \\
Déploiement & Docker Compose & une seule commande, identique sur un PC et sur un serveur \\
Graphiques & Recharts & courbes temps réel avec bande attendue \\
\bottomrule
\end{tabularx}
\caption{Choix technologiques}
\label{tab:choix}
\end{table}

\begin{aretenir}[Pourquoi pas PostgreSQL ou Kafka ?]
SentinelOps vise un seul hôte et une installation en une commande. Un serveur de base de
données ou un bus de messages aurait ajouté des conteneurs à configurer, ce que le projet
cherche justement à éviter. Le passage à PostgreSQL reste possible grâce aux migrations SQL
standard (voir les perspectives).
\end{aretenir}

\section{Extraits de code significatifs}

Le dépôt est organisé en quatre parties : \code{agent/} (Python), \code{server/} et \code{ui/} (Node.js et React), \code{shared/migrations/} (schéma SQL commun), plus le lanceur \code{scripts/sen.py}, le fichier \code{docker-compose.yml} et le pipeline \code{.github/workflows/}.

\subsection{Cycle de vie des alertes}

L'extrait~\ref{lst:alertmanager} montre le c\oe ur de l'automate : une condition reste en
attente tant que sa durée \code{for} n'est pas atteinte, puis l'alerte est ouverte. L'index
unique protège contre une ouverture concurrente.

\begin{lstlisting}[language=Python,caption={Passage de PENDING à FIRING (\texttt{alert\_manager.py})},label={lst:alertmanager}]
p = self.pending.get(key)
if p is None:
    p = self.pending[key] = Pending(since_ms=now_ms, hits=0)
p.hits += 1

if now_ms - p.since_ms < sig.for_seconds * 1000:
    stats["pending"] += 1
    continue                      # not long enough yet: nothing written

# PENDING -> FIRING (alert + outbox row in one transaction)
try:
    alert["id"] = await self.db.open_alert(alert)
except sqlite3.IntegrityError:
    continue                      # unique index: already open elsewhere
del self.pending[key]
\end{lstlisting}

\subsection{Score z robuste}

\begin{lstlisting}[language=Python,caption={Calcul du score z robuste (\texttt{anomaly\_detector.py})},label={lst:zscore}]
mean = statistics.median(samples)                          # robust centre
stdev = 1.4826 * statistics.median(abs(x - mean) for x in samples)
effective = max(stdev, self.min_noise_floor)               # floor = 2
z = (value - mean) / effective
is_outlier = (abs(z) > self.k_threshold                    # k = 3
              and abs(value - mean) >= self.min_abs_delta) # >= 8 units
\end{lstlisting}

\subsection{API fermée par défaut}

\begin{lstlisting}[language=JavaScript,caption={Intergiciel d'authentification (\texttt{middleware/auth.js})},label={lst:auth}]
const PUBLIC = new Set(['POST /auth/login', 'GET /health']);

export function requireAuth(req, res, next) {
  const key = `${req.method} ${req.path}`;
  if (PUBLIC.has(key)) return next();

  const user = authService.validateToken(bearerToken(req));
  if (!user) return res.status(401).json({ error: 'Authentication required' });

  if (user.must_change_password && !ONBOARDING.has(key))
    return res.status(403).json({ code: 'PASSWORD_CHANGE_REQUIRED' });
  req.user = user;
  next();
}
\end{lstlisting}

Une nouvelle route est donc protégée sans qu'on ait à y penser : il faudrait l'ajouter
explicitement à la liste \code{PUBLIC} pour la rendre accessible sans session.

La migration du schéma (verrou \code{BEGIN IMMEDIATE}) et le travailleur d'outbox (nouvel essai à 15, 30, 60 puis 120 s) suivent la conception du chapitre~\ref{chap:conception}.

\section{Lanceur en une commande}

Le lanceur \code{sen} (script Python sans dépendance, avec des raccourcis \code{sen} et
\code{sen.cmd}) masque Docker Compose. La commande \code{sen run} vérifie Docker, génère un mot
de passe administrateur au premier lancement, construit et démarre les conteneurs, puis attend
que \code{/api/health} réponde \textbf{et} que l'agent ait envoyé son premier signal de vie avant
d'afficher \og Ready \fg{} avec l'adresse et les identifiants (figure~\ref{fig:sen-run}).

\begin{table}[htbp]
\centering
\small
\begin{tabularx}{\linewidth}{L{4.4cm}Y}
\toprule
\textbf{Commande} & \textbf{Effet} \\
\midrule
\code{sen run} & construit, démarre et attend que tout soit prêt \\
\code{sen run --public --port 8080} & expose le dashboard sur le réseau, sur un autre port \\
\code{sen status} & état des conteneurs, de l'API et de l'agent \\
\code{sen logs [agent|dashboard]} & journaux en continu \\
\code{sen password} & affiche le mot de passe administrateur initial \\
\code{sen update}, \code{sen restart}, \code{sen stop} & mise à jour, redémarrage, arrêt \\
\bottomrule
\end{tabularx}
\caption{Commandes du lanceur \texttt{sen}}
\label{tab:sen}
\end{table}

\begin{figure}[htbp]
\centering
\souscapture[0.55\linewidth]{cap-18-sen-run.png}{\texttt{sen run} jusqu'à \og Ready \fg{}}\hfill
\souscapture[0.42\linewidth]{cap-19-sen-status.png}{\texttt{sen status}}
\caption{Le lanceur en action}
\label{fig:sen-run}
\end{figure}

\section{Interfaces de l'application}

Au premier lancement, l'administrateur se connecte avec le mot de passe généré puis doit en
choisir un nouveau ; tant que ce n'est pas fait, l'API refuse toute autre action. La vue
d'ensemble affiche ensuite l'état de l'hôte et de chaque conteneur, les incidents ouverts et
des courbes mises à jour en temps réel par WebSocket (figure~\ref{fig:ui-principal}).

\begin{figure}[htbp]
\centering
\souscapture[0.7\linewidth]{cap-02-login.png}{Connexion}\\[0.6em]
\souscapture{cap-04-vue-ensemble-normal.png}{Vue d'ensemble, état normal}\hfill
\souscapture{cap-05-vue-ensemble-incident.png}{Vue d'ensemble pendant un incident}
\caption{Authentification et vue d'ensemble}
\label{fig:ui-principal}
\end{figure}

Un clic sur un incident ouvre un panneau qui répond aux questions de l'opérateur : pourquoi
l'alerte s'est déclenchée (preuves mesurées), ce qui s'est passé, quels processus consommaient
les ressources, et un brouillon de compte rendu (\emph{post-mortem}). Les courbes affichent la
bande attendue calculée par le détecteur d'anomalies (figure~\ref{fig:ui-incident}).

\begin{figure}[htbp]
\centering
\souscapture{cap-06-detail-pourquoi.png}{Onglet \og Pourquoi \fg{}}\hfill
\souscapture{cap-08-detail-processus.png}{Instantané des processus}\\[0.6em]
\souscapture[0.7\linewidth]{cap-10-courbe-anomalie.png}{Courbe avec bande attendue et point anormal}
\caption{Analyse d'un incident}
\label{fig:ui-incident}
\end{figure}

Les règles, les canaux de notification et les heures calmes se gèrent depuis l'interface, ainsi
que les journaux et le redémarrage des conteneurs. L'interface existe aussi en thème clair et
s'adapte aux écrans de téléphone. Les notifications Slack signalent le déclenchement puis la
résolution (figure~\ref{fig:ui-gestion}).

\begin{figure}[htbp]
\centering
\souscapture{cap-11-regles.png}{Règles d'alerte}\hfill
\souscapture{cap-12-parametres.png}{Paramètres}\\[0.6em]
\souscapture{cap-13-journaux-conteneur.png}{Journaux d'un conteneur}\hfill
\souscapture{cap-16-slack-declenchement.png}{Notification Slack}
\caption{Gestion et notifications}
\label{fig:ui-gestion}
\end{figure}

\begin{conclusionchapitre}
Les fonctionnalités prévues au chapitre~\ref{chap:besoins} sont réalisées et accessibles par
une seule commande. Le chapitre suivant décrit comment elles sont sécurisées, testées et
déployées.
\end{conclusionchapitre}
